\subsection{Solubility from restricted coprime commutators: Problem 19.56}
\label{cand:19.56}
\problemstatement{19.56}

For finite \(G\), let
\[
 X(G)=\{[x,y]: |x|,\ |y|\text{ are powers of distinct primes}\}.
\]
Suppose \(|ab|\geq|a||b|\) for \(a,b\in X(G)\) of
coprime orders. Then \(G\) is soluble. In fact it
suffices to impose the inequality when the outputs
\(a,b\) also have prime-power orders.
The prime-power restrictions on the \emph{inputs}
are the additional issue here, compared with nearby
commutator criteria \cite{monakhov-derived,bastos-monetta}.

First, a primary element \(a\in X(G)\) normalizing a
subgroup \(R\) of different prime-power order must
centralize \(R\). For \(z\in R\), the element
\(b=[a,z]\) lies in \(R\cap X(G)\), while
\(ab=a^z\). If \(b\ne1\), the hypothesis gives
\(|a|=|ab|\geq|a||b|>|a|\), a contradiction.
This adapts the normalizer argument in
\cite{monakhov-derived}, with the required input
orders checked explicitly.

For \(P\in\operatorname{Syl}_p(G)\), let \(O^p(G)\)
be the \(p\)-residual, the intersection of normal
subgroups with \(p\)-group quotient. The hyperfocal
theorem in the finite-group form of
\cite[Lemma 2.2]{bcglo} gives
\[
 P\cap O^p(G)=
 \langle [x,g]:x\in Q\leq P,\ g\in N_G(Q),\
                                      (|g|,p)=1\rangle.
\]
Decompose \(g\) into its commuting primary
components inside \(\langle g\rangle\).
They still normalize \(Q\), and
\([x,uv]=[x,v][x,u]^v\) expresses each generator
as a product of elements of \(P\cap X(G)\).
Conversely, every \(X(G)\)-value vanishes in
a \(p\)-group quotient because at least one input
does. Hence
\[
 P\cap O^p(G)=\langle P\cap X(G)\rangle.
\]
For perfect \(G\), its \(p\)-group quotient is
trivial, so these primary \(X(G)\)-values generate
every Sylow subgroup.

Take a nonsoluble counterexample of least order.
The condition passes to subgroups, so every proper
subgroup is soluble and \(G=G'\).
Its soluble radical \(R\) is contained in every
maximal subgroup: otherwise \(G=RM\) makes
\(G/R\) soluble. Thus \(R\leq\Phi(G)\).
The Frattini subgroup of a finite group is
nilpotent: the Frattini argument for any of
its Sylow subgroups gives
\(G=\Phi(G)N_G(P)\), forcing \(N_G(P)=G\).
Thus \(R=\Phi(G)\), and \(G/R\) is minimal
nonabelian simple.

For a Sylow subgroup \(R_p\) of \(\Phi(G)\),
the normalizer observation makes every primary
\(X(G)\)-value of another prime centralize
\(R_p\). The generation result makes its
centralizer contain all the corresponding
Sylow subgroups of \(G\). That normal
centralizer has \(p\)-power index, and
perfection forces it to be \(G\). It follows that
\[
 \Phi(G)=Z(G),\qquad S=G/Z(G)\text{ is minimal simple}.
\]
The order condition has not been passed to a
quotient at any point.

We use the standard local consequence of Thompson's
minimal-simple classification recorded in
\cite[Proposition 18]{bastos-monetta}: a finite
minimal simple group contains a nontrivial
elementary abelian \(2\)-subgroup \(V\) with
a nontrivial action by an odd-order subgroup
normalizing it. Choose an element of that subgroup
acting nontrivially and decompose it into primary
components; at least one component \(\bar y\),
of odd prime-power order, acts nontrivially.
This weaker primary-order choice suffices here.
The archive also supplies explicit local matrices
for all five minimal-simple families, giving an
odd-prime-order choice.

The preimage of \(V\) in \(G\) has central derived
subgroup, hence is nilpotent of class at most two.
Its unique Sylow \(2\)-subgroup \(P\) maps onto
\(V\). Choose a primary lift \(y\) of \(\bar y\)
by taking the appropriate primary component of
any lift; it normalizes \(P\). For some \(u\in P\),
\[
 a=[u,y]\in P\cap X(G)
\]
has nonidentity image \(\bar a\in V\), an involution.

The Baer--Suzuki theorem, in its usual \(p\)-element
form recalled in \cite{guralnick-malle-baer}, gives
a conjugate \(\bar a^{\bar g}\) such that
\(\langle\bar a,\bar a^{\bar g}\rangle\) is not a
\(2\)-group, since \(O_2(S)=1\).
Two involutions generate a dihedral group.
Their product has order divisible by an odd
prime \(r\), and its appropriate power \(c\)
has order \(r\) and is inverted by \(\bar a\).

The preimage of \(\langle c\rangle\) is abelian,
as its quotient modulo the central kernel is
cyclic. Its unique Sylow \(r\)-subgroup \(R_r\)
maps onto \(\langle c\rangle\) and is normalized
by \(a\), in the preimage of the dihedral group.
The normalizer observation therefore makes \(a\)
centralize \(R_r\). In the quotient this says
that \(\bar a\) centralizes \(c\), contradicting
its inversion action on an element of odd order.
This proves the theorem.

The hyperfocal theorem, minimal-simple local
structure, and Baer--Suzuki theorem are imported
dependencies. The normalizer and central-extension
strategy has precedents in the cited works;
the required replacement is generation and
witness construction with coprime primary
inputs. Source: \artifact{research/19.56-proof.md}.
